\section{Identification with the exact density}
\label{subsec:rgnl-current-step}

The three terms in \eqref{eq:rgnl-full-step} have different origins:
the normalized incoming activity, the local determinant difference,
and the fluctuation logarithm. We first bound the determinant term,
then identify their evaluated sum with the exact block density.

The reblocked determinant term belongs to $K_{k+1}^x$.
Indeed the scaled arguments lie in its inverse domain,
and each summand is analytic and reads only fine
bonds meeting $X$. These scale to bonds meeting $Y$.
There are no legs. The lattice symmetry permutes the
reblocking summands; local gauge shifts of the two
slots become the corresponding real shifts on these
fine read bonds, where determinant locality and gauge
invariance apply. Evenness of $E_k^d$ gives charge
conjugation, and its conjugation identity gives reality.
The pointwise bound
\exteq{82}{FermionBlockManuscript} and
\exteq{83}{NormalizedContraction} show
\begin{align}
 \|\mathcal L\delta E_k^{\mathrm{det}}\|_{k+1}^x
    &\le8c_{\mathrm{rb}}L^3C_dM^3e_k^{1/4}
      \le e_k^{1/4-2\varepsilon}
      \le e_{k+1}^{1/4-2\varepsilon}.
                                      \label{eq:rgnl-reblocked-determinant}
\end{align}
Here $L(\kappa-\kappa_0-1)\ge\kappa$ controls the
remaining exponential weight, and the fixed charge
condition is
$8c_{\mathrm{rb}}L^3C_dM^3e_k^{2\varepsilon}\le1$.
Restriction to $\mathscr C=0$ gives the same result
for $\mathcal LE_k^{\mathrm{det}}$.

\begin{proof}[Proof of \cref{thm:rgnl-current-step}]
Use the notation of \cref{prop:rg-step-density-identity},
with $A=\mathscr A'_L$, $B=A+\mathscr C$ and
$Z=Z(1)$ the full-strength photon fluctuation.
The primitive families satisfy their exact identities
at full strength:
$H_k(1,B)=H_k(B)$, $\Gamma_k(1,B)=\Gamma_k(B)$, and
$\sum_XM_k(X,B)=b_k-D_k(B)$ for $k\ge1$.
The induced fields obey
\[
 \psi^{0,\#}=H_k^\#(B)\Psi^c,\qquad
 H_k^\#(B)\Psi^f=\psi^{0,\#}+W_k^\#(1,B).
\]
Define three evaluation homomorphisms. Let $\operatorname{ev}$
send old generators to $\psi^{0,\#}$ and fix $W$;
let $\operatorname{ev}_1$ send them to $H_k^\#(B)\Psi^f$;
for $k\ge1$, let $\operatorname{ev}_f$ additionally send
$\Psi,\bar\Psi$ to $\Psi^f,\bar\Psi^f$.
At zero take $\operatorname{ev}_f=\operatorname{ev}_1$.
These maps exist since their images are odd and have
the required antiperiodicity. They commute with the
finite coefficient integrals; $\operatorname{ev}$ also
commutes with Berezin integration, since it fixes $W$.

On the first contours $\mathcal C_rf=f(1)-f(0)$.
For $G_1$, summing the local primitives gives the
Dirac difference
$-\langle\bar\Psi,(D_k(B+Z)-D_k(B))\Psi\rangle_1$.
For $G_2,G_3$, evaluating
$\psi^\#+J_k^\#(1,B,Z)\Psi$ gives
$H_k^\#(B+Z)\Psi^f$.
The evaluations of their contour differences therefore give
\begin{equation}\label{eq:rgnl-first-contour-evaluation}
 \operatorname{ev}_f\sum_X\sum_{i\le3}G_i(1,X)
  =\mathrm{Act}_k^x[m',\mathscr E^r](A+Z,\mathscr C;\Psi^f)
       -\mathrm{Act}_k^x[m',\mathscr E^r](A,\mathscr C;\Psi^f).
\end{equation}
At zero the explicit first vertex gives this same
identity with the chamber pairing: $D_0(B)=D_B+\mu$
and the mass form has no connection dependence.

The reconstruction left inverse gives, on every local
first expression,
$\operatorname{ev}_1\operatorname{sub}_{Y,B}
=\operatorname{ev}_f$.
This follows on the $\Psi$ generators from $T_k(B)H_k(B)=1$
and on the old generators from their definition.
Adding $W_k^\#(1,B)$ to $\psi^\#$ in a polynomial and then
applying $\operatorname{ev}$ gives
$\operatorname{ev}_1$ applied to that polynomial.
Apply these identities to
\eqref{eq:rgnl-two-localization-sums}.
The evaluated $G_{\le3}^{\,W}$ sum is
\eqref{eq:rgnl-first-contour-evaluation}.
The evaluated polynomial contours $G_4,G_5$ give
\[
 I_k^x[m',\mathscr E^r](A,\mathscr C;\Psi^f)
                      -I_k^x[m',\mathscr E^r](A,\mathscr C;\Psi^c).
\]
Consequently
\[
 \operatorname{ev}\sum_YF(Y)=V_k^x,\qquad
 \Xi_k^x
     =\exp\left\{\sum_X
                  \operatorname{ev}E_k^\#(m,\mathscr E)(X;A,\mathscr C)\right\}.
\]
The second equality is
\eqref{eq:rgnl-log-identity} evaluated by
$\operatorname{ev}$; its evaluation commutes with the
integrals as just noted.

The other remaining factor is the determinant quotient.
Equation~\exteq{78}{FermionBlockManuscript} gives
\[
 \frac{\delta Z_k(B)}{\delta Z_k(0)}
      =\exp\left\{\sum_X\delta E_k^{\mathrm{det}}(X;A,\mathscr C)
                               +\sum_XE_k^{\mathrm{det}}(X,\mathscr C)\right\}.
\]
On inserted fields adjustment changes no value:
$\operatorname{ev}R\mathscr E(X)
=\operatorname{ev}(R\mathscr E)^{\mathrm{nat}}(X)$.
The equality is the telescoping reconstruction identity
for a small polymer and the definition for a large one.
Finally $\operatorname{ev}=\operatorname{ev}'S_L$,
where $\operatorname{ev}'$ inserts the outgoing fields.
Each fine polymer has exactly one outgoing reblocking
polymer. Thus the three activity sums combine into the
outgoing $\mathscr E^+$, and the two scalar external sums
combine into \eqref{eq:rgnl-external-scalar-map}.
Insert these identities in
\eqref{eq:rg-step-density-identity}. All exponents are
even and commute, so their sums give
\eqref{eq:rgnl-exact-step}, with exactly its previously
derived mass, vacuum and determinant normalization.

Integrate over the outgoing generators and apply the exact
marginal \eqref{eq:rg-step-block-marginal}. This gives
\eqref{eq:rgnl-integrated-step}.
The integrand is admissible throughout the photon box:
Theorem~\extref{B}{MaxwellBlockManuscript} places
$A_k^{\min}(A')+\mathrm{Sq}_k\zeta$ in
$\mathrm{Dom}_k$, and its reconstructed connection is
an affine, continuous function of $\zeta$.
The density coefficients are analytic in this connection
on its operator and activity domains, hence continuous
and integrable on the compact box.
The external scalar factor is fixed during this integral.
No condition on the harmonic part of $A'$ enters.

The three terms of \eqref{eq:rgnl-full-step} belong to
$K_{k+1}^x$ by Proposition~\extref{6.1}{NormalizedContraction},
\eqref{eq:rgnl-reblocked-determinant}, and
\cref{prop:rgnl-scaled-correction}.
Their respective norms are bounded by the three terms
in \eqref{eq:rgnl-activity-step-bound}.
Equation~\exteq{46}{NormalizedContraction} bounds the
first term with the stronger exponent $3\varepsilon/2$.
The mass and vacuum extraction bounds give
\eqref{eq:rgnl-mass-step-bound} and
\eqref{eq:rgnl-vacuum-step-bound}, since
\[
 0\le\delta\varepsilon_k^0
        \le12\exp(-p_{0,k}^2/2)\le e_k^7.
\]
At zero all incoming extracted terms vanish.
The exact algebra uses all $192n_0$ pairs with the
weight $1/48$ and integrates to the $4n_1$-pair
outgoing density. There is no removed multiplicity
or additional vacuum normalization. The displayed
zero-level formulas and bounds follow.
\end{proof}

The restriction to zero external connection commutes
with all three terms of \eqref{eq:rgnl-full-step}.
The determinant difference there becomes
$E_k^d(X,\mathscr A)-E_k^d(X,0)$, and the external scalar
map preserves $\mathfrak P_k(\cdot,0)=0$.
Thus \eqref{eq:rgnl-exact-step} and
\eqref{eq:rgnl-integrated-step} also give the step without
a slot, with incoming activity $\mathscr E|_0$.

The displayed norm estimates depend on the running charge
and the fixed constants, and on no side or harmonic part.
They were proved for any primitive system with the stated hypotheses on
the present torus. Identifying the outputs from different
systems or different volumes requires the separate
comparison arguments; it does not follow from these bounds.


